% Phase I data section, publication form. Every quantity here was read from the file that
% holds it; the paths are named so each can be checked.
\section{Data}\label{sec:data}

\subsection{Discharge observations and the training network}

The hourly training network comprises \num{9181} gauges compiled into a single archive of
\num{403248} hourly steps. The record begins at 00:00 UTC on 1 January 1980 and ends at
23:00 UTC on 31 December 2025, covering \num{16802} days, that is \num{46} complete
calendar years with no gap in the time axis. Six national and regional collections contribute, in descending order of
size: CAMELS-H for the conterminous United States (\num{5278} gauges), the Australian
Bureau of Meteorology network (\num{1840}), LamaH-CE for central Europe (\num{835}), a
Japanese national collection (\num{690}), a German collection (\num{463}), and LamaH-Ice for
Iceland (\num{75}).

The prediction target is specific discharge $q$, expressed in \si{\mm\per\hour} at hourly
resolution and \si{\mm\per\day} at daily resolution. Normalising by catchment area in this
way makes gauges of different size directly comparable and lets one model serve all of them.

Gauges whose observed hourly series has a standard deviation below \num{e-3} are excluded
from every reported metric, since a variance ratio and a correlation are both undefined for a
constant series. That filter removes \num{136} gauges. The \num{8843} that remain constitute
the target-domain population for all target-domain results in this report.

The geography of this network is a limitation of the study and belongs here rather than in a
closing caveat. The network contains no gauge in Africa, none in South America and none in
mainland Asia. Approximately four fifths of its gauges lie between \SI{30}{\degree} and
\SI{60}{\degree} north. The word \emph{global} in this report describes the intended scope of
the model. It does not describe the observational coverage available to train it. Section
\ref{sec:africa-design} exists to address that gap directly.

\subsection{Meteorological forcing and catchment attributes}

Three hourly meteorological variables drive the model: precipitation, air temperature, and
potential evapotranspiration. Precipitation is taken from MSWEP version 3.16. Potential
evapotranspiration is computed by the Penman method rather than by a temperature-only
approximation, since the latter degrades most in exactly the arid and tropical settings the
African test emphasises.

Forty-two static catchment attributes are available per gauge. Seventeen are withheld from
the model. Sixteen of these are socio-economic proxies: gridded population counts and
densities, gross domestic product at two resolutions, the Human Development Index, and the
pedigree fields that accompany them. They are withheld because a model may exploit them to
identify a region rather than to characterise a catchment, which would inflate apparent
generalisation. The seventeenth withheld field is a duplicate elevation product already
represented by another attribute. Twenty-five attributes remain in use.

\subsection{Sample construction}

A sample is one target hour together with the meteorological history preceding it. The
model receives two views of that history. The daily branch receives the \num{365} days
immediately before the target day. The hourly branch receives the final \SI{336}{\hour}. The
forcing at the target hour itself is withheld from the input, so that forcing up to hour
$t-1$ predicts discharge at hour $t$.

Target hours are sampled with a stride of \num{24} at hour 23 of each day, which places each
sample's final \num{24} hourly outputs on one calendar day and allows consecutive samples to
be joined into a continuous hourly series without overlap or gaps. That property is what
makes the within-day diagnostics of Section~\ref{sec:falsify} possible. The construction
yields \num{73808280} samples over the full network.

Training and validation are separated in time within each gauge rather than by holding out
whole gauges. The first \SI{70}{\percent} of a gauge's valid samples train and the final
\SI{30}{\percent} validate, giving \num{51665658} training samples and \num{22142622}
validation samples. A temporal split of this form ensures that no validation sample precedes
a training sample at the same gauge, so the evaluation cannot benefit from having seen the
future of its own series.

\subsection{Construction of the daily branch}

Prepared training batches were available and were used in a first set of experiments. In
those batches the daily branch is not a daily series. It is a \num{1000}-step subsample of
the hourly record drawn under a power law, dense in the hours immediately before the target
and increasingly sparse further back, spanning \num{8760} hours in total.

That subsample cannot supply a daily mean, and the reason is quantitative rather than a
matter of preference. Distributing the \num{1000} sampled points across the \num{365} days
they cover gives the following: \num{8} days retain all \num{24} hours, \num{176} days retain
exactly one point, \num{7} days retain none, and the remainder lie between. For
precipitation, where the majority of hours are zero, a single instantaneous sample is a poor
estimator of the daily mean, and it is biased in an uncontrolled direction depending on
whether the sampled hour happened to be wet.

The windows were therefore rebuilt from the raw hourly source, computing the daily branch as
$\bar{x}_d = \frac{1}{24}\sum_{h=1}^{24} x_{d,h}$ over each of the \num{365} preceding days.
Both constructions were carried through to completion and both are reported, in
Appendix~\ref{app:config}. The prepared archive is retained in the rebuilt path for its
normalisation constants and feature names alone, so that scores from the two constructions
remain directly comparable.

\subsection{Cross-validation design}

Gauges are partitioned into five folds, and each fold serves once as the held-out set while
the remaining four fifths supply training. Every gauge therefore takes exactly one turn as a
held-out gauge, and the target-domain metrics reported here cover the entire network rather
than a single fifth of it. No result depends on which fifth happened to be drawn.

Two partitioning schemes are compared. The first assigns gauges to folds at random, which
leaves each held-out gauge a trainable neighbour at a median distance of \SI{10.4}{\km}. The
second first clusters gauge coordinates into \num{120} groups by $k$-means on the unit sphere
and then assigns whole clusters to folds, raising the median distance to the nearest
trainable neighbour to \SI{94.9}{\km}. All other settings are identical between the two. The
difference between them isolates how much apparent zero-shot skill derives from spatial
proximity to training data rather than from generalisation across catchment characteristics.

\subsection{The African evaluation set}

\num{294} African catchments constitute the external evaluation set. They are drawn from
three archives: GRDC-Caravan (\num{126} catchments), the Global Runoff Data Centre
(\num{119}), and the restricted African Database of Hydrometric Indices (\num{49}).
\num{282} of them satisfy a minimum of \num{100} valid days within the evaluation period and
are scored, with a median of \num{1090} validation days each. They are small to medium
catchments, with a median area of \SI{781}{\km\squared} and a range from \num{19} to
\num{9835}.

The size of this set reflects data availability rather than selection. The daily discharge
database contains \num{302} African records with a time series; the \num{294} used here are
\SI{97}{\percent} of that total. A further \num{1283} African entries carry station metadata
without any accompanying series and cannot be evaluated by any method. Africa is genuinely
sparse in gauge coverage, and that sparsity is the condition the external test is designed to
confront rather than an artefact of how the set was assembled.

The particular \num{294} are the evaluation set of a published continent-holdout experiment,
adopted without modification so that results reported here rest against that study's baseline
on identical catchments. Whatever selection criteria that study applied are inherited, which
is a limitation of the comparison and is noted as such. No African catchment appears anywhere
in pretraining, and none carries hourly discharge observations.

Meteorological forcing over these catchments is ERA5-Land at \SI{0.1}{\degree} hourly
resolution, aggregated to catchment means by weighting each grid cell by its fractional area
of overlap with the catchment polygon. It runs from 01:00 UTC on 1 January 1980 to 23:00 UTC
on 31 December 1995, \num{140255} hourly steps. The African evaluation period is bounded by
this forcing rather than by the discharge records, several of which extend later. ERA5-Land total runoff over
the same catchments is retained separately as a process-based baseline. Because that variable
accumulates from 00 UTC and resets daily, hourly values are recovered by differencing within
each accumulation day, and the resulting hour-ending stamps are shifted to hour-beginning so
that they align with the model's output convention.
